\documentclass[11pt,a4paper]{article}

\usepackage[utf8]{inputenc}
\usepackage[
  left=1.05in,
  right=1.05in,
  top=1in,
  bottom=1in,
  headheight=14pt
]{geometry}
\usepackage{amsmath,amssymb,amsthm,mathtools}
\usepackage{bm}
\usepackage{enumitem}
\usepackage{booktabs,array}
\usepackage{microtype}
\usepackage{xcolor}
\usepackage[breakable,skins]{tcolorbox}
\usepackage[hidelinks]{hyperref}
\usepackage{fancyhdr}

\linespread{1.25}
\setlength{\parindent}{0pt}
\setlength{\parskip}{1em plus 0.25em minus 0.1em}
\setlength{\abovedisplayskip}{1.2em plus 0.3em minus 0.2em}
\setlength{\belowdisplayskip}{1.2em plus 0.3em minus 0.2em}
\setlength{\abovedisplayshortskip}{0.9em plus 0.2em}
\setlength{\belowdisplayshortskip}{1em plus 0.2em}
\setlength{\jot}{9pt}
\setlength{\emergencystretch}{2em}
\raggedbottom

\setlist[itemize]{leftmargin=*,itemsep=6pt,topsep=6pt,parsep=3pt}
\setlist[enumerate]{leftmargin=*,itemsep=6pt,topsep=6pt,parsep=3pt}

% Plain black frame for the question statement.
\newtcolorbox{qbox}{
  enhanced,
  colback=white,colframe=black,
  boxrule=0.6pt,sharp corners,
  left=10pt,right=10pt,top=4pt,bottom=6pt,
  before skip=12pt,after skip=12pt
}

\newenvironment{solution}
  {\par\noindent\textbf{Solution.}\par\nobreak}
  {\par\hfill\(\square\)\par}

\newcommand{\step}[1]{\par\addvspace{0.5em}\noindent\textbf{#1}\par\nobreak}
\newcommand{\mk}[1]{\hfill\textit{[#1 marks]}}

\newcommand{\pd}[2]{\frac{\partial #1}{\partial #2}}
\newcommand{\pdd}[2]{\frac{\partial^2 #1}{\partial #2^2}}
\newcommand{\norm}[1]{\lVert#1\rVert}

\pagestyle{fancy}
\fancyhf{}
\fancyhead[L]{MTL712}
\fancyhead[R]{Major Examinations: Solutions}
\fancyfoot[C]{\thepage}
\renewcommand{\headrulewidth}{0.4pt}
\renewcommand{\footrulewidth}{0pt}

\title{\textbf{MTL712: Computational Methods for Differential\\Equations}\\[6pt]
\large Previous-Year Major Papers with Complete Solutions}
\author{\normalsize Papers: November 2016, May 2023 and 2025--26 Semester I}
\date{}

\begin{document}

\maketitle
\thispagestyle{plain}

Questions are transcribed from the supplied scans.  Each question is
stated in a box, followed by a complete solution.  Where a question is
ambiguous or contains a misprint, the interpretation used is stated
before the solution.

{\linespread{1.0}\selectfont\setlength{\parskip}{0pt}\tableofcontents}
\clearpage

%#####################################################################
\section{Major Examination, November 2016}
%#####################################################################

Duration 2 hours, 50 marks.

%---------------------------------------------------------------------
\subsection{RK2 applied to a linear system}

\begin{qbox}
Find the condition on the step size \(h\) if the two-stage
second-order RK method applied to the system \(\mathbf y'=-A\mathbf y\),
where
\[
  A=\begin{bmatrix}2&0\\0&6\end{bmatrix},
\]
is absolutely stable.  What is the corresponding result for
\(A=\begin{bmatrix}20&0\\0&60\end{bmatrix}\)?\mk{7}
\end{qbox}

\begin{solution}
\step{Step 1: the stability function of any two-stage RK2 method}
A two-stage explicit RK method is
\[
  k_1=f(x_n,y_n),\qquad
  k_2=f(x_n+c_2h,\ y_n+c_2hk_1),\qquad
  y_{n+1}=y_n+h(b_1k_1+b_2k_2),
\]
and it has order \(2\) exactly when \(b_1+b_2=1\) and \(b_2c_2=\tfrac12\).
For the test equation \(y'=\lambda y\), with \(z=h\lambda\),
\[
  k_1=\lambda y_n,\qquad k_2=\lambda(1+c_2z)y_n,
\]
so
\[
  y_{n+1}=\bigl[1+(b_1+b_2)z+b_2c_2z^2\bigr]y_n
         =\Bigl(1+z+\frac{z^2}2\Bigr)y_n .
\]
Hence every two-stage second-order RK method has the same stability
function
\[
  R(z)=1+z+\frac{z^2}{2}.
\]

\step{Step 2: the real interval of absolute stability}
For real \(z\),
\[
  R(z)=\frac{(1+z)^2+1}{2}\ \ge\ \frac12>0,
\]
so \(|R(z)|<1\iff R(z)<1\iff z+\frac{z^2}2<0\iff z(2+z)<0\iff -2<z<0\).
The interval of absolute stability is \((-2,0)\).

\step{Step 3: decouple the system}
Since \(A\) is diagonal, \(\mathbf y'=-A\mathbf y\) is two independent
scalar equations, and the RK method acts on each component separately.
The components have \(\lambda_1=-2\) and \(\lambda_2=-6\).  The method is
absolutely stable for the system if and only if it is absolutely stable
for each component:
\[
  -2<-2h<0\quad\text{and}\quad -2<-6h<0
  \iff 0<h<\tfrac13 .
\]
The stiffer component (\(\lambda=-6\)) decides the bound.

\step{Step 4: the second matrix}
Now \(\lambda_1=-20\), \(\lambda_2=-60\), and
\[
  -2<-60h<0\iff 0<h<\tfrac1{30}.
\]

\step{Comment}
Multiplying \(A\) by \(10\) multiplies the eigenvalues by \(10\) and
divides the admissible step by \(10\): for an explicit method the step
size is limited by the largest eigenvalue, \(h<2/\lambda_{\max}\), even
though the corresponding solution component decays fastest.  This is
the difficulty with stiff systems.
\end{solution}

%---------------------------------------------------------------------
\subsection{Solving the implicit equation of an LMM}

\begin{qbox}
Consider a \(k\)-step LMM
\[
  y_{n+k}+\alpha_{k-1}y_{n+k-1}+\cdots+\alpha_0y_n
  =h\bigl(\beta_kf_{n+k}+\beta_{k-1}f_{n+k-1}+\cdots+\beta_0f_n\bigr)
\]
used to solve \(y'(x)=f(x,y)\), \(y(x_0)=y_0\), where \(f(x,y)\) is
non-linear.  Compare the fixed point iteration method, the
predictor--corrector method and the Newton--Raphson method for solving
the above problem.\mk{7}
\end{qbox}

\begin{solution}
If \(\beta_k\ne0\) the method is implicit.  At each step the unknown
\(y_{n+k}\) satisfies the non-linear equation
\begin{equation}\label{eq:impl}
  y_{n+k}=h\beta_kf(x_{n+k},y_{n+k})+g_n,
  \qquad
  g_n=h\sum_{j=0}^{k-1}\beta_jf_{n+j}-\sum_{j=0}^{k-1}\alpha_jy_{n+j},
\end{equation}
where \(g_n\) is known.  Let \(L\) be the Lipschitz constant of \(f\)
with respect to \(y\).

\step{(a) Fixed point iteration}
\[
  y_{n+k}^{[s+1]}=h\beta_kf\bigl(x_{n+k},y_{n+k}^{[s]}\bigr)+g_n,
  \qquad s=0,1,2,\dots
\]
The map \(y\mapsto h\beta_kf(x_{n+k},y)+g_n\) has Lipschitz constant
\(h|\beta_k|L\), so by the contraction mapping theorem the iteration
converges to the unique solution of \eqref{eq:impl} for any starting
value if
\[
  h<\frac1{|\beta_k|L}.
\]
\begin{itemize}
  \item Convergence is linear, with rate \(h|\beta_k|L\).
  \item Cost per iteration: one evaluation of \(f\); no derivatives.
  \item Drawback: the condition \(h|\beta_k|L<1\) restricts \(h\).  For
    stiff problems \(L\) is very large, so this destroys the advantage
    of using an implicit method.
\end{itemize}

\step{(b) Predictor--corrector method}
An explicit LMM (the predictor) gives a starting value
\(y_{n+k}^{[0]}\), and the implicit method (the corrector) is applied a
fixed number \(m\) of times, as in fixed point iteration:
\[
  \text{P(EC)}^m\text{E}:\quad
  \text{Predict},\ (\text{Evaluate }f,\ \text{Correct})^m,\ \text{Evaluate}.
\]
\begin{itemize}
  \item Cheap and fully explicit: \(m+1\) evaluations of \(f\) per step
    (typically \(m=1\) or \(2\)), no linear systems.
  \item The iteration is \emph{not} carried to convergence, so the
    result is a new method.  If the predictor has order \(p^*\ge p\)
    (the order of the corrector), the pair has order \(p\) and the
    principal error of the corrector; if \(p^*<p\), the order can drop.
  \item The difference between predicted and corrected values gives a
    cheap error estimate (Milne's device), useful for step control.
  \item The stability region is that of the PC pair, which is much
    smaller than that of the corrector and behaves like an explicit
    method.  So it is not suitable for stiff problems.
\end{itemize}

\step{(c) Newton--Raphson method}
Write \eqref{eq:impl} as \(F(y)=y-h\beta_kf(x_{n+k},y)-g_n=0\).  Then
\[
  \Bigl[I-h\beta_k\frac{\partial f}{\partial y}\bigl(x_{n+k},y^{[s]}\bigr)\Bigr]
  \bigl(y^{[s+1]}-y^{[s]}\bigr)=-F\bigl(y^{[s]}\bigr).
\]
\begin{itemize}
  \item Converges quadratically from a good starting value (for example
    from an explicit predictor).
  \item No restriction of the type \(h|\beta_k|L<1\): the method keeps
    the full stability of the implicit formula.  This makes it the
    method of choice for stiff problems.
  \item Cost per iteration: evaluation of the Jacobian
    \(\partial f/\partial y\) and solution of a linear system.  This can
    be reduced by \emph{modified Newton} (keeping the Jacobian fixed for
    several iterations or steps), at the price of linear convergence.
\end{itemize}

\step{Summary}
\begin{center}
\renewcommand{\arraystretch}{1.3}
\begin{tabular}{@{}lccc@{}}
\toprule
 & Fixed point & PC (P(EC)\(^m\)E) & Newton--Raphson\\
\midrule
Convergence & linear, needs \(h|\beta_k|L<1\) & fixed \(m\) corrections & quadratic (local)\\
Cost per step & \(f\) evaluations & \(m+1\) evaluations & Jacobian + linear solve\\
Step restriction & yes & yes (explicit-type) & no\\
Stiff problems & unsuitable & unsuitable & suitable\\
\bottomrule
\end{tabular}
\end{center}
For non-stiff problems the predictor--corrector approach is the most
economical; for stiff problems Newton's method (or modified Newton) is
needed.
\end{solution}

%---------------------------------------------------------------------
\subsection{Von Neumann criterion in the \(\ell_{2,\Delta x}\) norm}

\begin{qbox}
Prove that the difference scheme
\[
  u^{n+1}=Qu^n
\]
is stable with respect to the \(\ell_{2,\Delta x}\) norm if and only if
there exist positive constants \(\Delta t_0\) and \(\Delta x_0\) and
non-negative constants \(\beta\) and \(K\) so that
\[
  |\rho(\xi)|^{n+1}\le Ke^{\beta(n+1)\Delta t}
\]
for \(0<\Delta t\le\Delta t_0\), \(0<\Delta x\le\Delta x_0\) and all
\(\xi\in[-\pi,\pi]\).\mk{7}
\end{qbox}

\begin{solution}
\step{Setting}
\(Q\) is a constant-coefficient difference operator on the whole line,
\[
  (Qu)_k=\sum_{j=-m_1}^{m_2}q_ju_{k+j},
\]
where the \(q_j\) may depend on \(\Delta t\) and \(\Delta x\).  The norm
is \(\norm{u}_{2,\Delta x}^2=\sum_k|u_k|^2\Delta x\).  The scheme is
\emph{stable} if there are \(\Delta t_0,\Delta x_0>0\) and
\(K,\beta\ge0\) such that
\[
  \norm{u^{n+1}}_{2,\Delta x}\le Ke^{\beta(n+1)\Delta t}\norm{u^0}_{2,\Delta x}
\]
for all \(u^0\), all \(n\ge0\), \(0<\Delta t\le\Delta t_0\) and
\(0<\Delta x\le\Delta x_0\).

\step{Discrete Fourier transform}
For \(u\in\ell_2\) define
\[
  \hat u(\xi)=\frac1{\sqrt{2\pi}}\sum_k e^{-ik\xi}u_k,\qquad \xi\in[-\pi,\pi].
\]
Parseval's identity gives \(\norm{\hat u}_{L_2[-\pi,\pi]}=\norm{u}_{\ell_2}\),
and hence \(\norm{u}_{2,\Delta x}=\sqrt{\Delta x}\,\norm{\hat u}_{L_2}\).
Shifting by \(j\) multiplies the transform by \(e^{ij\xi}\), so
\[
  \widehat{u^{n+1}}(\xi)=\rho(\xi)\,\widehat{u^n}(\xi),
  \qquad
  \rho(\xi)=\sum_jq_je^{ij\xi},
\]
and therefore \(\widehat{u^{n+1}}=\rho^{n+1}\,\widehat{u^0}\).  The
symbol \(\rho\) is a trigonometric polynomial, hence continuous in
\(\xi\).

\step{(\(\Leftarrow\)) The condition implies stability}
If \(|\rho(\xi)|^{n+1}\le Ke^{\beta(n+1)\Delta t}\) for all \(\xi\), then
\begin{align*}
  \norm{u^{n+1}}_{2,\Delta x}^2
  &=\Delta x\int_{-\pi}^{\pi}|\rho(\xi)|^{2(n+1)}\,|\widehat{u^0}(\xi)|^2\,d\xi\\
  &\le K^2e^{2\beta(n+1)\Delta t}\,\Delta x\int_{-\pi}^{\pi}|\widehat{u^0}(\xi)|^2\,d\xi
   =K^2e^{2\beta(n+1)\Delta t}\norm{u^0}_{2,\Delta x}^2 .
\end{align*}
Taking square roots gives stability with the same \(K\) and \(\beta\).

\step{(\(\Rightarrow\)) Stability implies the condition}
Suppose the scheme is stable with constants \(\Delta t_0,\Delta x_0,K,\beta\),
and suppose, for contradiction, that for some admissible
\(\Delta t,\Delta x\), some \(n\) and some \(\xi_0\),
\[
  |\rho(\xi_0)|^{n+1}>Ke^{\beta(n+1)\Delta t}.
\]
By continuity of \(\rho\) there is an interval \(I\ni\xi_0\) of positive
length on which \(|\rho(\xi)|^{n+1}>Ke^{\beta(n+1)\Delta t}\).  Choose
\(u^0\) to be the sequence whose transform is the indicator function of
\(I\) (it lies in \(\ell_2\) by Parseval).  Then
\[
  \norm{u^{n+1}}_{2,\Delta x}^2
  =\Delta x\int_I|\rho(\xi)|^{2(n+1)}\,d\xi
  >K^2e^{2\beta(n+1)\Delta t}\,\Delta x\int_I d\xi
  =K^2e^{2\beta(n+1)\Delta t}\norm{u^0}_{2,\Delta x}^2,
\]
contradicting stability.  Hence the inequality holds, with the same
constants, for all admissible \(\Delta t,\Delta x\) and all \(\xi\).
\end{solution}

\textit{Remark.}  In practice the condition is checked through the
sufficient form \(|\rho(\xi)|\le1+C\Delta t\): then
\(|\rho|^{n+1}\le(1+C\Delta t)^{n+1}\le e^{C(n+1)\Delta t}\), so
\(K=1\), \(\beta=C\).

%---------------------------------------------------------------------
\subsection{Crank--Nicolson for \(v_t=v_{xx}-v\)}

\begin{qbox}
Consider the PDE \(v_t=v_{xx}-v\), \(t>0\), \(x\in(-\infty,\infty)\).
Analyse the stability of the Crank--Nicolson scheme for this PDE.\mk{7}
\end{qbox}

\begin{solution}
\step{The scheme}
Crank--Nicolson averages the whole right-hand side between the levels
\(n\) and \(n+1\).  With \(r=\Delta t/\Delta x^2\) and
\(\delta^2u_k=u_{k+1}-2u_k+u_{k-1}\),
\[
  \frac{u_k^{n+1}-u_k^n}{\Delta t}
  =\frac12\Bigl[\frac{\delta^2u_k^{n+1}}{\Delta x^2}-u_k^{n+1}\Bigr]
   +\frac12\Bigl[\frac{\delta^2u_k^{n}}{\Delta x^2}-u_k^{n}\Bigr],
\]
that is,
\[
  \Bigl(1+\frac{\Delta t}2\Bigr)u_k^{n+1}-\frac r2\,\delta^2u_k^{n+1}
  =\Bigl(1-\frac{\Delta t}2\Bigr)u_k^{n}+\frac r2\,\delta^2u_k^{n}.
\]

\step{The symbol}
Substitute \(u_k^n=\rho^ne^{ik\xi}\).  Since
\(\delta^2e^{ik\xi}=(2\cos\xi-2)e^{ik\xi}=-4\sin^2\frac\xi2\,e^{ik\xi}\),
\[
  \rho(\xi)
  =\frac{1-\frac{\Delta t}2-2r\sin^2\frac\xi2}
        {1+\frac{\Delta t}2+2r\sin^2\frac\xi2}
  =\frac{1-a}{1+a},
  \qquad
  a=2r\sin^2\frac\xi2+\frac{\Delta t}2 .
\]

\step{Stability}
Since \(a>0\) for every \(\xi\) and every \(\Delta t>0\),
\[
  |1-a|<1+a\quad\Longrightarrow\quad|\rho(\xi)|<1 .
\]
Hence \(|\rho(\xi)|^{n+1}\le1\), and by the criterion of the previous
question (with \(K=1\), \(\beta=0\)) the scheme is stable for all
\(\Delta t\) and \(\Delta x\): it is \emph{unconditionally stable}.

This agrees with the PDE, whose Fourier modes decay like
\(e^{-(1+\omega^2)t}\).  (For large \(r\), \(\rho\) is close to \(-1\)
for high frequencies, so these modes decay slowly and oscillate in
sign, as for Crank--Nicolson applied to the heat equation.)
\end{solution}

%---------------------------------------------------------------------
\subsection{Consistency with a derivative boundary condition}

\begin{qbox}
Discuss the consistency of the difference scheme
\[
  u_k^{n+1}=(1-2r)u_k^n+r\bigl(u_{k+1}^n+u_{k-1}^n\bigr),\quad
  k=1,2,\dots,M-1,\qquad r=\frac{\nu\Delta t}{\Delta x^2},
\]
\[
  u_M^{n+1}=0,\qquad u_0^n=u_1^n,
\]
to the initial-boundary value problem
\[
  v_t=\nu v_{xx},\ x\in(0,1),\ t>0;\qquad
  v(x,0)=f(x);\qquad v(1,t)=0;\qquad v_x(0,t)=0 .
\]\mk{7}
\end{qbox}

\begin{solution}
A scheme for an IBVP is consistent if the difference equation and each
difference boundary condition are consistent with the corresponding
equation and boundary condition of the problem.  Let \(v\) be a smooth
exact solution and \(v_k^n=v(k\Delta x,n\Delta t)\).

\step{Interior points \(k=1,\dots,M-1\)}
Written in the form that approximates the PDE,
\[
  \tau_k^n=\frac{v_k^{n+1}-v_k^n}{\Delta t}
  -\nu\,\frac{v_{k+1}^n-2v_k^n+v_{k-1}^n}{\Delta x^2}.
\]
By Taylor expansion about \((k\Delta x,n\Delta t)\),
\[
  \tau_k^n=(v_t-\nu v_{xx})+\frac{\Delta t}2v_{tt}
  -\frac{\nu\Delta x^2}{12}v_{xxxx}+\cdots
  =O(\Delta t)+O(\Delta x^2),
\]
since \(v_t=\nu v_{xx}\).  So \(\tau_k^n\to0\) as
\(\Delta t,\Delta x\to0\): the interior scheme is consistent, of order
\((1,2)\).

\step{The boundary \(x=1\)}
\(u_M^{n+1}=0\) imposes \(v(1,t)=0\) exactly; the truncation error is
zero.

\step{The boundary \(x=0\)}
The condition \(u_0^n=u_1^n\) is the approximation
\[
  \frac{u_1^n-u_0^n}{\Delta x}=0\qquad\text{of}\qquad v_x(0,t)=0 .
\]
Its truncation error is
\[
  \frac{v(\Delta x,t)-v(0,t)}{\Delta x}-v_x(0,t)
  =\frac{\Delta x}{2}v_{xx}(0,t)+O(\Delta x^2)=O(\Delta x),
\]
which tends to zero.  So the boundary condition is consistent, but only
\emph{first order} accurate (it is a one-sided difference).

\step{Conclusion}
The scheme is consistent with the IBVP.  The interior equation is
\(O(\Delta t)+O(\Delta x^2)\) accurate, but the Neumann boundary
condition is only \(O(\Delta x)\), which in general reduces the overall
accuracy to first order in \(\Delta x\).  A second-order treatment uses
a fictitious point: \(u_{-1}^n=u_1^n\) (central difference), which gives
\(u_0^{n+1}=(1-2r)u_0^n+2ru_1^n\) at \(k=0\).
\end{solution}

%---------------------------------------------------------------------
\subsection{FTCS and FTBS for \(v_t+av_x=\nu v_{xx}\)}

\begin{qbox}
Consider the IVP \(v_t+av_x=\nu v_{xx}\), \(v(x,0)=f(x)\).
\begin{enumerate}[label=(\roman*)]
  \item Analyse the stability of the FTCS scheme for the above IVP given
    that \(r^2\ge R^2/4\), where \(r=\nu\Delta t/\Delta x^2\) and
    \(R=a\Delta t/\Delta x\).\mk{5}
  \item Analyse the stability of the FTBS scheme when \(\nu=0\) and
    \(a>0\).\mk{5}
  \item Analyse the stability of the FTCS scheme when \(\nu=0\) and
    \(a>0\).\mk{4}
  \item From (iii), what can you say about the stability of the FTCS
    scheme for the IBVP \(v_t+av_x=0\), \(v(x,0)=f(x)\), \(a>0\),
    \(v(0,t)=v(1,t)=0\)?  Give reasons.\mk{1}
\end{enumerate}
\end{qbox}

\begin{solution}
Throughout, \(s=\sin\frac\xi2\) and \(\xi\in[-\pi,\pi]\).

\step{(i) FTCS for convection--diffusion}
The scheme is
\[
  u_k^{n+1}=u_k^n-\frac R2\bigl(u_{k+1}^n-u_{k-1}^n\bigr)
  +r\bigl(u_{k+1}^n-2u_k^n+u_{k-1}^n\bigr).
\]
Its symbol is
\[
  \rho(\xi)=1-2r(1-\cos\xi)-iR\sin\xi=1-4rs^2-iR\sin\xi .
\]
Using \(\sin^2\xi=4s^2(1-s^2)\),
\begin{align*}
  |\rho|^2&=(1-4rs^2)^2+4R^2s^2(1-s^2)\\
  &=1-4s^2\bigl[(2r-R^2)-(4r^2-R^2)s^2\bigr].
\end{align*}
So \(|\rho|\le1\) for all \(\xi\) if and only if
\(g(s^2)=(2r-R^2)-(4r^2-R^2)s^2\ge0\) for all \(s^2\in[0,1]\).  \(g\) is
linear in \(s^2\), so this holds if and only if
\[
  g(0)=2r-R^2\ge0\qquad\text{and}\qquad g(1)=2r-4r^2\ge0,
\]
that is,
\[
  R^2\le2r\le1 .
\]
Under the given assumption \(R^2\le4r^2\): if \(r\le\frac12\) then
\(R^2\le4r^2=2r\cdot2r\le2r\), so the first condition holds
automatically.  Hence, given \(r^2\ge R^2/4\),
\[
  \text{FTCS is stable}\iff 0<r\le\tfrac12 .
\]

\step{(ii) FTBS for \(v_t+av_x=0\), \(a>0\)}
The scheme is \(u_k^{n+1}=u_k^n-R\bigl(u_k^n-u_{k-1}^n\bigr)\), with
symbol
\[
  \rho(\xi)=1-R+Re^{-i\xi}.
\]
Then
\[
  |\rho|^2=(1-R+R\cos\xi)^2+R^2\sin^2\xi=1-4R(1-R)\sin^2\frac\xi2 .
\]
So \(|\rho|\le1\) for all \(\xi\iff R(1-R)\ge0\iff0\le R\le1\).  Since
\(a>0\), FTBS is stable if and only if
\[
  0<R=\frac{a\Delta t}{\Delta x}\le1
\]
(the CFL condition).  Geometrically, \(\rho\) is a circle of radius \(R\)
centred at \(1-R\), which lies inside the unit disc exactly when
\(0\le R\le1\).

\step{(iii) FTCS for \(v_t+av_x=0\)}
With \(\nu=0\) the symbol is \(\rho(\xi)=1-iR\sin\xi\), so
\[
  |\rho(\xi)|^2=1+R^2\sin^2\xi>1\qquad\text{whenever }\sin\xi\ne0 .
\]
For the usual refinement with \(R\) fixed, \(\max_\xi|\rho|=\sqrt{1+R^2}\)
is a constant greater than \(1\), independent of \(\Delta t\), so
\(|\rho|^{n+1}=(1+R^2)^{(n+1)/2}\) cannot be bounded by
\(Ke^{\beta(n+1)\Delta t}\).  FTCS is \emph{unconditionally unstable}.
(Only under the impractical restriction \(\Delta t\le c\,\Delta x^2\) do
we get \(R^2=a^2\Delta t^2/\Delta x^2\le a^2c\,\Delta t\), hence
\(|\rho|\le1+\frac12a^2c\,\Delta t\), which is stable in the weak sense.)

\step{(iv) The IBVP}
The FTCS scheme is still unstable.  On the grid \(x_k=k\Delta x\),
\(M\Delta x=1\), with \(u_0=u_M=0\), the scheme is
\(\mathbf u^{n+1}=Q\mathbf u^n\) with
\[
  Q=I-\frac R2S,\qquad S=\mathrm{tridiag}(-1,0,1),
\]
\(S\) skew-symmetric.  By the tridiagonal eigenvalue formula the
eigenvalues of \(S\) are \(2i\cos\frac{j\pi}M\), so those of \(Q\) are
\[
  \lambda_j=1-iR\cos\frac{j\pi}M,
  \qquad |\lambda_j|^2=1+R^2\cos^2\frac{j\pi}M,\quad j=1,\dots,M-1 .
\]
\(Q\) is normal (identity plus a skew-symmetric matrix), so
\(\norm{Q^n}_2=\max_j|\lambda_j|^n=\bigl(1+R^2\cos^2\frac\pi M\bigr)^{n/2}\),
which grows geometrically in \(n\) for fixed \(R\).  This is the
instability found in (iii).  In addition, for \(a>0\) the
characteristics carry information to the right, so a condition may be
prescribed only at the inflow boundary \(x=0\); the condition
\(v(1,t)=0\) over-determines the problem.
\end{solution}

%#####################################################################
\section{Final Examination, May 2023}
%#####################################################################

\textit{Note.}  This paper (a second-semester offering) covers numerical
linear algebra, interpolation, quadrature and optimisation rather than
differential equations; these topics are not in the lecture notes of
this course.  Full solutions are given anyway.  Reasons must be given
for all answers.

%---------------------------------------------------------------------
\subsection{Gradient descent with exact line search}

\begin{qbox}
Consider the gradient descent method applied to a convex function
\(f:\mathbb R^n\to\mathbb R\), where at each step we move along the
negative gradient direction to a point which minimises \(f\) along this
line.  Let \(x_k\) and \(x_{k+1}\) be two consecutive points generated
by this method.  Prove that \(\nabla f(x_k)^T\nabla f(x_{k+1})=0\).\mk{3}
\end{qbox}

\begin{solution}
Assume \(f\) is differentiable and the line minimum exists.  Let
\(g=\nabla f(x_k)\).  If \(g=0\), the result is trivial.  Otherwise
\[
  x_{k+1}=x_k-t_kg,\qquad t_k\in\arg\min_{t\ge0}\varphi(t),
  \qquad \varphi(t)=f(x_k-tg).
\]
By the chain rule, \(\varphi'(t)=-\nabla f(x_k-tg)^Tg\).  In particular
\(\varphi'(0)=-\norm g^2<0\), so \(\varphi\) decreases near \(t=0\) and
the minimiser satisfies \(t_k>0\).  An interior minimiser of a
differentiable function is a stationary point, so
\[
  0=\varphi'(t_k)=-\nabla f(x_k-t_kg)^Tg=-\nabla f(x_{k+1})^T\nabla f(x_k).
\]
(Convexity of \(f\) makes \(\varphi\) convex, so any stationary point
of \(\varphi\) is indeed its minimiser.)  Hence consecutive gradients
are orthogonal: the method zig-zags.
\end{solution}

%---------------------------------------------------------------------
\subsection{Newton's method for \(f(x)=\norm x_2^3\)}

\begin{qbox}
Determine the rate of convergence of Newton's method (with step size
\(\eta=1\)) for minimising \(f(x)=\norm x_2^3\), \(x\in\mathbb R^n\).
You may use: if \(u\) is a unit vector,
\((I+uu^T)^{-1}=I-\frac12uu^T\).\mk{4}
\end{qbox}

\begin{solution}
For \(x\ne0\) write \(\norm x=\norm x_2\) and \(u=x/\norm x\).  Then
\[
  \nabla f(x)=3\norm x\,x,
  \qquad
  \nabla^2f(x)=3\Bigl(\norm x\,I+\frac{xx^T}{\norm x}\Bigr)
  =3\norm x\,(I+uu^T).
\]
By the given formula,
\[
  [\nabla^2f(x)]^{-1}=\frac1{3\norm x}\Bigl(I-\frac12uu^T\Bigr).
\]
The Newton step is
\[
  x_{k+1}=x_k-[\nabla^2f(x_k)]^{-1}\nabla f(x_k)
  =x_k-\Bigl(I-\frac12uu^T\Bigr)x_k
  =\frac12\,u\,(u^Tx_k)=\frac12x_k ,
\]
since \(u^Tx_k=\norm{x_k}\).  Hence
\[
  x_k=2^{-k}x_0\to0=x^*,\qquad \frac{\norm{x_{k+1}-x^*}}{\norm{x_k-x^*}}=\frac12 .
\]
Newton's method converges only \emph{linearly}, with rate \(\frac12\),
not quadratically.  The reason is that the Hessian at the minimiser is
singular (\(\nabla^2f(0)=0\)), so the hypothesis of the quadratic
convergence theorem fails.
\end{solution}

%---------------------------------------------------------------------
\subsection{Convexity of a maximum}

\begin{qbox}
\begin{enumerate}[label=(\roman*)]
  \item Let \(f_1,\dots,f_m\) be convex functions with a common domain.
    Prove that \(f(x)=\max_{i}f_i(x)\) is also convex.\mk{3}
  \item Given \(x\in\mathbb R^n\), define \(f(x)\) as the sum of the
    \(k\) largest coordinates of \(x\), where \(k\) is a fixed integer
    between \(0\) and \(n\).  Prove that \(f\) is convex.\mk{2}
\end{enumerate}
\end{qbox}

\begin{solution}
\step{(i)}
The common domain \(D\) is convex.  For \(x,y\in D\) and
\(\theta\in[0,1]\), for every \(i\),
\[
  f_i(\theta x+(1-\theta)y)\le\theta f_i(x)+(1-\theta)f_i(y)
  \le\theta f(x)+(1-\theta)f(y).
\]
Taking the maximum over \(i\) on the left gives
\(f(\theta x+(1-\theta)y)\le\theta f(x)+(1-\theta)f(y)\).

\step{(ii)}
For any set \(S\subseteq\{1,\dots,n\}\) with \(|S|=k\), the sum
\(\sum_{i\in S}x_i\) is at most the sum of the \(k\) largest
coordinates, with equality when \(S\) indexes them.  Hence
\[
  f(x)=\max_{|S|=k}\ \sum_{i\in S}x_i ,
\]
a maximum of finitely many (\(\binom nk\)) linear, hence convex,
functions.  By (i), \(f\) is convex.  (For \(k=0\), \(f\equiv0\).)
\end{solution}

%---------------------------------------------------------------------
\subsection{Interpolation error for \(\cos x\)}

\begin{qbox}
Suppose \(\cos(x)\) is to be approximated by an interpolating
polynomial of degree \(n+1\), using \(n\) equally spaced points in the
interval \([0,1]\).  How accurate is the approximation (give bounds, no
exact answer is needed)?\mk{2}
\end{qbox}

\begin{solution}
\textit{Interpretation.}  A polynomial of degree \(n+1\) needs \(n+2\)
interpolation points, so the statement is inconsistent.  We give the
bound for \(N+1\) equally spaced nodes in \([0,1]\) (degree \(N\)) and
then specialise.

\step{Error formula}
For \(x\in[0,1]\) there is \(\eta\in(0,1)\) with
\[
  \cos x-p_N(x)=\frac{f^{(N+1)}(\eta)}{(N+1)!}\prod_{i=0}^N(x-x_i).
\]
Every derivative of \(\cos\) is \(\pm\sin\) or \(\pm\cos\), so
\(|f^{(N+1)}|\le1\).

\step{Crude bound}
Each factor satisfies \(|x-x_i|\le1\), so
\[
  |\cos x-p_N(x)|\le\frac1{(N+1)!}.
\]

\step{Sharper bound for equal spacing}
With spacing \(h=1/N\), \(\bigl|\prod_{i=0}^N(x-x_i)\bigr|\le\frac14h^{N+1}N!\),
so
\[
  |\cos x-p_N(x)|\le\frac{h^{N+1}}{4(N+1)}=\frac1{4(N+1)N^{N+1}}.
\]

\step{Specialising}
For degree \(n+1\) (\(N=n+1\), i.e.\ \(n+2\) points):
\[
  |\cos x-p(x)|\le\frac1{(n+2)!}
  \quad\text{and}\quad
  |\cos x-p(x)|\le\frac1{4(n+2)(n+1)^{n+2}} .
\]
If instead \(n\) points are used literally (degree \(n-1\)), the bounds
are \(1/n!\) and \(1/\bigl(4n(n-1)^n\bigr)\).  Either way the error
decreases faster than geometrically.  For example, a polynomial of
degree \(5\) on \(6\) equally spaced points (\(h=\frac15\)) has error at
most \(1/(4\cdot6\cdot5^6)\approx2.7\times10^{-6}\) (the crude bound
gives \(1/720\approx1.4\times10^{-3}\)).
\end{solution}

%---------------------------------------------------------------------
\subsection{A quadrature formula with derivatives}

\begin{qbox}
Determine a formula of the form
\[
  \int_a^bf(x)\,dx=Af(c)+Bf'(a)+Cf''(b),
\]
where \(c\) is the mid-point of \([a,b]\).  The formula should be exact
for polynomials of as high degree as possible.\mk{3}
\end{qbox}

\begin{solution}
Shift so that the midpoint is \(0\): let \(H=\frac{b-a}2\), so the
interval is \([-H,H]\) and \(c=0\) (the coefficients do not change
under translation).  Impose exactness for \(1,x,x^2\):
\begin{align*}
  f=1:&\quad 2H=A,\\
  f=x:&\quad 0=A\cdot0+B\cdot1+C\cdot0\ \Rightarrow\ B=0,\\
  f=x^2:&\quad \tfrac23H^3=A\cdot0+B\cdot(-2H)+2C\ \Rightarrow\ C=\tfrac13H^3 .
\end{align*}
Check \(f=x^3\): the integral is \(0\), but the right-hand side is
\(A\cdot0+B\cdot3H^2+C\cdot6H=2H^4\ne0\).  So the maximal degree of
exactness is \(2\), and
\[
  \boxed{\int_a^bf(x)\,dx\approx(b-a)\,f\Bigl(\frac{a+b}2\Bigr)
  +\frac{(b-a)^3}{24}\,f''(b)},
\]
exact for all polynomials of degree \(\le2\) (and not for \(x^3\)).
\end{solution}

%---------------------------------------------------------------------
\subsection{An LU factorisation}

\begin{qbox}
Compute the LU factorisation without pivoting of
\(\begin{pmatrix}1&2\\2&5\end{pmatrix}\).\mk{2}
\end{qbox}

\begin{solution}
Eliminate with multiplier \(l_{21}=2/1=2\): row 2 \(-\,2\times\) row 1
gives \((0,\ 5-4)=(0,1)\).  Hence
\[
  \begin{pmatrix}1&2\\2&5\end{pmatrix}
  =\begin{pmatrix}1&0\\2&1\end{pmatrix}
   \begin{pmatrix}1&2\\0&1\end{pmatrix}.
\]
(Check: \(2\cdot2+1\cdot1=5\).)
\end{solution}

%---------------------------------------------------------------------
\subsection{A fixed point iteration}

\begin{qbox}
Consider the fixed point iteration \(x_{k+1}=(\alpha+1)x_k-x_k^2\),
\(k=0,1,\dots\), where \(\alpha\) is a real number satisfying
\(\frac12\le\alpha\le1\).  Show that the iteration converges for any
initial guess \(x_0\) with \(\alpha-\frac15\le x_0\le\alpha+\frac15\).
Assuming that the iteration converges, find the values of \(\alpha\) for
which it converges quadratically.  Argue why it will converge
quadratically for this range of the starting point.\mk{5}
\end{qbox}

\begin{solution}
\step{Fixed points}
\(g(x)=(\alpha+1)x-x^2\).  Then \(g(x)=x\iff x(\alpha-x)=0\), so the
fixed points are \(0\) and \(\alpha\).  Since \(g'(0)=1+\alpha>1\), the
fixed point \(0\) is repelling; the iteration converges to \(\alpha\).

\step{Error recursion}
Let \(e_k=x_k-\alpha\).  Then
\[
  e_{k+1}=g(\alpha+e_k)-\alpha
  =(\alpha+1)(\alpha+e_k)-(\alpha+e_k)^2-\alpha
  =e_k\,(1-\alpha-e_k).
\]

\step{Convergence for \(|e_0|\le\frac15\)}
If \(|e_k|\le\frac15\) and \(\frac12\le\alpha\le1\), then
\[
  -\tfrac15\le1-\alpha-e_k\le\tfrac12+\tfrac15=\tfrac7{10},
\]
so \(|1-\alpha-e_k|\le\frac7{10}\) and \(|e_{k+1}|\le\frac7{10}|e_k|\).
By induction, \(|e_k|\le\frac15\) for all \(k\) (the iterates stay in
\([\alpha-\frac15,\alpha+\frac15]\)) and
\[
  |e_k|\le\Bigl(\tfrac7{10}\Bigr)^k|e_0|\to0 .
\]
So the iteration converges to \(\alpha\) for every such \(x_0\).

\step{Quadratic convergence}
From the recursion,
\(e_{k+1}=(1-\alpha)e_k-e_k^2\).  If \(\alpha<1\), then
\(e_{k+1}/e_k\to1-\alpha=g'(\alpha)\ne0\): the convergence is linear
with rate \(1-\alpha\).  Convergence is quadratic exactly when
\(g'(\alpha)=1-\alpha=0\), i.e.
\[
  \alpha=1 .
\]
Then \(e_{k+1}=-e_k^2\), so \(|e_{k+1}|=|e_k|^2\) exactly.  For
\(|e_0|\le\frac15\) this gives \(|e_k|=|e_0|^{2^k}\le5^{-2^k}\to0\): the
iteration converges, and quadratically, with asymptotic constant
\(|e_{k+1}|/|e_k|^2=1\).
\end{solution}

%---------------------------------------------------------------------
\subsection{A Householder reflection}

\begin{qbox}
Determine the Householder transformation that annihilates all but the
first entry of the vector \((1,1,1,1)^T\), i.e.\ find the Householder
matrix \(H\) such that \(H(1,1,1,1)^T=(\alpha,0,0,0)^T\), where
\(\alpha\) is a negative number.  Clearly write down the vector
corresponding to the Householder matrix.\mk{3}
\end{qbox}

\begin{solution}
Let \(x=(1,1,1,1)^T\).  A reflection preserves length, so
\(|\alpha|=\norm x_2=2\), and \(\alpha=-2\) (the sign opposite to
\(x_1\), which also avoids cancellation).  The Householder vector is
\[
  v=x-\alpha e_1=x+2e_1=(3,1,1,1)^T,\qquad v^Tv=12,
\]
and
\[
  H=I-\frac{2vv^T}{v^Tv}=I-\frac16vv^T
  =\frac16\begin{pmatrix}
    -3&-3&-3&-3\\
    -3&5&-1&-1\\
    -3&-1&5&-1\\
    -3&-1&-1&5
  \end{pmatrix}.
\]
(Equivalently, the unit vector is \(u=v/\norm v=(3,1,1,1)^T/(2\sqrt3)\)
and \(H=I-2uu^T\).)  Check: the first row of \(Hx\) is
\(\frac16(-12)=-2\), and the second row is
\(\frac16(-3+5-1-1)=0\); similarly for rows 3 and 4.  So
\(Hx=(-2,0,0,0)^T\).
\end{solution}

%---------------------------------------------------------------------
\subsection{Power iteration}

\begin{qbox}
Consider the matrix \(A=\begin{pmatrix}1&4\\1&1\end{pmatrix}\).  To which
eigenvector will the power iteration starting with \((1,1)^T\)
converge?\mk{3}
\end{qbox}

\begin{solution}
The characteristic equation is \((1-\lambda)^2-4=0\), so
\(\lambda_1=3\), \(\lambda_2=-1\).  Eigenvectors:
\[
  (A-3I)v=0:\ v_1=(2,1)^T;\qquad (A+I)v=0:\ v_2=(2,-1)^T .
\]
Expand the starting vector:
\((1,1)^T=c_1(2,1)^T+c_2(2,-1)^T\) gives \(2c_1+2c_2=1\),
\(c_1-c_2=1\), so \(c_1=\frac34\), \(c_2=-\frac14\).  Then
\[
  A^k\begin{pmatrix}1\\1\end{pmatrix}
  =3^k\Bigl[\tfrac34v_1-\tfrac14\Bigl(-\tfrac13\Bigr)^kv_2\Bigr].
\]
Since \(c_1\ne0\) and \(|\lambda_1|=3>|\lambda_2|=1\), the normalised
iterates converge to the dominant eigenvector
\[
  \frac{(2,1)^T}{\sqrt5},
\]
corresponding to \(\lambda=3\), with error \(O(3^{-k})\).  (Because
\(\lambda_1>0\), the iterates converge without sign alternation.)
\end{solution}

%---------------------------------------------------------------------
\subsection{Backward stability of \(QA\)}

\begin{qbox}
Let \(Q\) be an \(n\times n\) unitary matrix (\(QQ^*=I\)).  Consider the
computational problem: given an \(n\times n\) matrix \(A\), output
\(QA\).  Assume that you are given stable algorithms for computing
addition and multiplication.  Prove that the usual algorithm for
computing \(QA\) is stable.  (The input is \(A\); \(Q\) is fixed.)\mk{5}
\end{qbox}

\begin{solution}
(The paper writes \(UU^*=I\); this means \(QQ^*=I\).)

\step{Floating point model}
Each operation satisfies \(\mathrm{fl}(a\circ b)=(a\circ b)(1+\delta)\)
with \(|\delta|\le\varepsilon\) (machine epsilon), for
\(\circ\in\{+,\times\}\).

\step{Error in one inner product}
The usual algorithm computes each entry
\((QA)_{ij}=\sum_{k=1}^nq_{ik}a_{kj}\) by multiplying and summing in
order.  Each product picks up one factor \((1+\delta)\), and each term
passes through at most \(n-1\) additions, each contributing a factor
\((1+\delta)\).  Hence
\[
  \mathrm{fl}\Bigl(\sum_kq_{ik}a_{kj}\Bigr)=\sum_kq_{ik}a_{kj}(1+\theta_k),
  \qquad |\theta_k|\le\gamma_n:=\frac{n\varepsilon}{1-n\varepsilon}
  =n\varepsilon+O(\varepsilon^2).
\]
So the computed product is \(\widetilde{QA}=QA+E\) with
\[
  |E_{ij}|\le\gamma_n\sum_k|q_{ik}||a_{kj}|,
  \qquad\text{i.e.}\qquad |E|\le\gamma_n|Q|\,|A|\ \text{(entrywise)}.
\]

\step{Backward error}
Define \(\delta A=Q^*E\).  Then
\[
  Q(A+\delta A)=QA+QQ^*E=QA+E=\widetilde{QA},
\]
so the computed result is the \emph{exact} product for the perturbed
input \(A+\delta A\).  Since \(Q^*\) is unitary,
\(\norm{\delta A}_F=\norm{E}_F\), and
\[
  \norm{E}_F\le\gamma_n\,\bigl\lVert|Q|\,|A|\bigr\rVert_F
  \le\gamma_n\norm{Q}_F\norm{A}_F=\gamma_n\sqrt n\,\norm A_F,
\]
using \(\norm{Q}_F^2=\mathrm{trace}(Q^*Q)=n\).  Therefore
\[
  \frac{\norm{\delta A}_F}{\norm A_F}\le\sqrt n\,\gamma_n=O(\varepsilon).
\]

\step{Conclusion}
The algorithm is \emph{backward stable}: it gives the exact answer for
an input with relative perturbation \(O(\varepsilon)\).  Backward
stability implies stability.  (Unitarity was essential: it made the
backward error \(Q^*E\) no larger than the forward error \(E\).)
\end{solution}

\clearpage
%#####################################################################
\section{End-Semester Examination, 2025--26 Semester I}
%#####################################################################

Duration 2 hours, 40 marks.  All notations are standard.

%---------------------------------------------------------------------
\subsection{A tridiagonal system solved as a difference equation}

\begin{qbox}
Suppose we want to solve the tridiagonal system \(Ax=b\), where
\[
  A=\begin{bmatrix}
    -3&1&&&\\
    2&-3&1&&\\
    &2&-3&1&\\
    &&\ddots&\ddots&\ddots\\
    &&2&-3&1\\
    &&&2&-3
  \end{bmatrix},
  \qquad
  b=\begin{bmatrix}1\\0\\0\\\vdots\\0\\0\end{bmatrix}.
\]
State the difference equation which replaces the matrix formulation of
the problem.  Solve the system by solving the difference equation.\mk{4}
\end{qbox}

\begin{solution}
Let \(A\) be \(N\times N\) and \(x=(x_1,\dots,x_N)^T\).

\step{Step 1: the difference equation}
Row \(i\) of \(Ax=b\) reads
\[
  2x_{i-1}-3x_i+x_{i+1}=b_i,\qquad i=1,\dots,N,
\]
where the first and last rows have no \(x_0\) and \(x_{N+1}\) terms.
The only non-zero entry of \(b\) is \(b_1=1\), in the first row
\(-3x_1+x_2=1\).  Introduce two fictitious values \(x_0\), \(x_{N+1}\)
and choose them so that every row has the same form: the first row is
\(2x_0-3x_1+x_2=0\) provided \(2x_0=-1\), and the last row is
\(2x_{N-1}-3x_N+x_{N+1}=0\) provided \(x_{N+1}=0\).  Hence \(Ax=b\) is
equivalent to the boundary value problem for a second-order
homogeneous linear difference equation:
\[
  \boxed{\;x_{i+1}-3x_i+2x_{i-1}=0,\quad i=1,2,\dots,N,
  \qquad x_0=-\tfrac12,\quad x_{N+1}=0.\;}
\]

\step{Step 2: the general solution}
In shift-operator form, \((E^2-3E+2)x_{i-1}=0\).  The characteristic
equation \(\xi^2-3\xi+2=(\xi-1)(\xi-2)=0\) has the distinct roots
\(\xi=1,2\), so
\[
  x_i=C_1+C_2\,2^i .
\]

\step{Step 3: the boundary conditions}
\[
  x_0=C_1+C_2=-\tfrac12,\qquad x_{N+1}=C_1+C_2\,2^{N+1}=0 .
\]
Subtracting, \(C_2(2^{N+1}-1)=\frac12\), so
\[
  C_2=\frac1{2(2^{N+1}-1)},\qquad
  C_1=-2^{N+1}C_2=-\frac{2^N}{2^{N+1}-1}.
\]
(The determinant \(2^{N+1}-1\ne0\), so the solution is unique, as it
must be since \(A\) is non-singular.)  Therefore
\[
  \boxed{\;x_i=\frac{2^{i-1}-2^N}{2^{N+1}-1},\qquad i=1,2,\dots,N.\;}
\]

\step{Check}
First row:
\(-3x_1+x_2=\dfrac{-3(1-2^N)+(2-2^N)}{2^{N+1}-1}=\dfrac{2^{N+1}-1}{2^{N+1}-1}=1\).
Last row:
\(2x_{N-1}-3x_N=\dfrac{(2^{N-1}-2^{N+1})-3(2^{N-1}-2^N)}{2^{N+1}-1}=0\).
For example, \(N=2\) gives \(x=(-\frac37,-\frac27)^T\), and indeed
\(-3(-\frac37)-\frac27=1\), \(2(-\frac37)-3(-\frac27)=0\).  All
\(x_i<0\), and \(|x_i|\) decreases from \(x_1\approx-\frac12\) to
\(x_N\approx-\frac14\) for large \(N\).
\end{solution}

%---------------------------------------------------------------------
\subsection{Classical RK4 for a rotating system}

\begin{qbox}
The classical fourth-order Runge--Kutta method is used to solve the
system
\[
  y'=-kz,\quad y(x_0)=y_0,\qquad z'=ky,\quad z(x_0)=z_0,
\]
where \(k>0\) and \(x,x_0,y,y_0,z\) and \(z_0\) are real.  The step
length \(h\) is supposed to be \(>0\).  Putting
\(y_n\approx y(x_0+nh)\) and \(z_n\approx z(x_0+nh)\), prove that
\[
  \begin{bmatrix}y_{n+1}\\ z_{n+1}\end{bmatrix}
  =A\begin{bmatrix}y_n\\ z_n\end{bmatrix},
\]
where \(A\) is a real \(2\times2\) matrix.  Find under what
condition(s) the solutions do not grow exponentially for increasing
values of \(n\).\mk{5}
\end{qbox}

\begin{solution}
\step{Step 1: vector form}
Let \(\mathbf w=(y,z)^T\).  The system is \(\mathbf w'=M\mathbf w\) with
\[
  M=kJ,\qquad J=\begin{bmatrix}0&-1\\1&0\end{bmatrix},\qquad
  J^2=-I,\quad J^3=-J,\quad J^4=I .
\]

\step{Step 2: RK4 on a linear constant-coefficient system}
The classical RK4 method for \(\mathbf w'=\mathbf F(x,\mathbf w)\) is
\[
  \mathbf w_{n+1}=\mathbf w_n+\frac h6\bigl(K_1+2K_2+2K_3+K_4\bigr),
\]
\begin{align*}
  K_1&=\mathbf F(x_n,\mathbf w_n),&
  K_2&=\mathbf F\bigl(x_n+\tfrac h2,\mathbf w_n+\tfrac h2K_1\bigr),\\
  K_3&=\mathbf F\bigl(x_n+\tfrac h2,\mathbf w_n+\tfrac h2K_2\bigr),&
  K_4&=\mathbf F(x_n+h,\mathbf w_n+hK_3).
\end{align*}
With \(\mathbf F=M\mathbf w\),
\begin{align*}
  K_1&=M\mathbf w_n,\\
  K_2&=M\bigl(\mathbf w_n+\tfrac h2M\mathbf w_n\bigr)
      =\bigl(M+\tfrac h2M^2\bigr)\mathbf w_n,\\
  K_3&=M\bigl(\mathbf w_n+\tfrac h2K_2\bigr)
      =\bigl(M+\tfrac h2M^2+\tfrac{h^2}4M^3\bigr)\mathbf w_n,\\
  K_4&=M\bigl(\mathbf w_n+hK_3\bigr)
      =\bigl(M+hM^2+\tfrac{h^2}2M^3+\tfrac{h^3}4M^4\bigr)\mathbf w_n .
\end{align*}
Hence
\[
  K_1+2K_2+2K_3+K_4
  =\bigl(6M+3hM^2+h^2M^3+\tfrac{h^3}4M^4\bigr)\mathbf w_n,
\]
and
\[
  \mathbf w_{n+1}
  =\Bigl(I+hM+\frac{(hM)^2}{2!}+\frac{(hM)^3}{3!}+\frac{(hM)^4}{4!}\Bigr)\mathbf w_n .
\]

\step{Step 3: the matrix \(A\)}
Put \(\omega=kh>0\).  Using \(J^2=-I\), \(J^3=-J\), \(J^4=I\),
\[
  A=\Bigl(1-\frac{\omega^2}2+\frac{\omega^4}{24}\Bigr)I
   +\Bigl(\omega-\frac{\omega^3}6\Bigr)J
  =\begin{bmatrix}a&-b\\ b&a\end{bmatrix},
  \qquad
  a=1-\frac{\omega^2}2+\frac{\omega^4}{24},\quad
  b=\omega-\frac{\omega^3}6 .
\]
This is a real \(2\times2\) matrix, which proves the first part.
Explicitly,
\[
  y_{n+1}=a\,y_n-b\,z_n,\qquad z_{n+1}=b\,y_n+a\,z_n .
\]

\step{Step 4: growth of the solution}
\(A^TA=(a^2+b^2)I\), so \(A=\sqrt{a^2+b^2}\,R\) with \(R\) a rotation,
and the eigenvalues of \(A\) are \(a\pm ib\), of modulus
\(\sqrt{a^2+b^2}\).  Therefore
\[
  y_n^2+z_n^2=\bigl(a^2+b^2\bigr)^n\bigl(y_0^2+z_0^2\bigr)
\]
exactly, for every \(n\).  Now
\begin{align*}
  a^2&=1-\omega^2+\frac{\omega^4}3-\frac{\omega^6}{24}+\frac{\omega^8}{576},\qquad
  b^2=\omega^2-\frac{\omega^4}3+\frac{\omega^6}{36},\\
  a^2+b^2&=1-\frac{\omega^6}{72}+\frac{\omega^8}{576}
  =1-\frac{\omega^6}{576}\bigl(8-\omega^2\bigr).
\end{align*}
The solutions do not grow exponentially if and only if
\(a^2+b^2\le1\), i.e.\ \(\omega^2\le8\):
\[
  \boxed{\;0<kh\le2\sqrt2,\qquad\text{i.e.}\qquad 0<h\le\frac{2\sqrt2}{k}\approx\frac{2.83}{k}.\;}
\]
\begin{itemize}
  \item \(kh<2\sqrt2\): \(a^2+b^2<1\), the numerical solution decays
    slowly (a small amplitude error).
  \item \(kh=2\sqrt2\): \(a^2+b^2=1\), \(A\) is a rotation and
    \(y_n^2+z_n^2\) is constant.
  \item \(kh>2\sqrt2\): \(y_n^2+z_n^2\) grows like
    \((a^2+b^2)^n\), exponentially.
\end{itemize}

\step{Comment}
The exact solution
\(y=y_0\cos k(x-x_0)-z_0\sin k(x-x_0)\),
\(z=y_0\sin k(x-x_0)+z_0\cos k(x-x_0)\) rotates with constant
\(y^2+z^2\).  The eigenvalues of \(M\) are \(\pm ik\), so the result is
the statement that the RK4 stability polynomial
\(R(\bar h)=1+\bar h+\frac{\bar h^2}2+\frac{\bar h^3}6+\frac{\bar h^4}{24}\)
satisfies \(|R(i\omega)|\le1\) exactly for \(|\omega|\le2\sqrt2\): the
interval of the imaginary axis inside the RK4 stability region.
\end{solution}

%---------------------------------------------------------------------
\subsection{Order of an LMM via \texorpdfstring{\(\rho(z)/\log z-\sigma(z)\)}{rho/log - sigma}}

\begin{qbox}
Consider the initial value problem \(y'=f(x,y)\), \(y(x_0)=y_0\).
Assume \(f(x,y)\) is a sufficiently smooth function.  Suppose \(p\) is
a positive integer.  Prove that the linear \(k\)-step method defined by
\(\sum_{j=0}^k\alpha_jy_{n+j}=h\sum_{j=0}^k\beta_jf_{n+j}\), with
stability polynomials \(\rho(z)\) and \(\sigma(z)\), is of order \(p\)
if and only if the function \(\varphi(z)\) defined by
\[
  \varphi(z)=\frac{\rho(z)}{\log z}-\sigma(z)
\]
has a zero of multiplicity \(p\) at \(z=1\).  (Note that the function
\(\log z\) appearing in the denominator is made single valued by
cutting the complex plane along the half line \(\operatorname{Re}z\le0\).)\mk{5}
\end{qbox}

\begin{solution}
\step{Definitions}
\(\rho(z)=\sum_{j=0}^k\alpha_jz^j\), \(\sigma(z)=\sum_{j=0}^k\beta_jz^j\).
The associated difference operator is
\[
  L[y(x);h]=\sum_{j=0}^k\bigl[\alpha_jy(x+jh)-h\beta_jy'(x+jh)\bigr].
\]
Expanding \(y(x+jh)\) and \(y'(x+jh)\) in Taylor series about \(x\),
\[
  L[y(x);h]=c_0y(x)+c_1hy'(x)+c_2h^2y''(x)+\cdots,
\]
\[
  c_0=\sum_j\alpha_j,\qquad
  c_q=\frac1{q!}\sum_jj^q\alpha_j-\frac1{(q-1)!}\sum_jj^{q-1}\beta_j\quad(q\ge1).
\]
The method has \emph{order \(p\)} if \(c_0=c_1=\cdots=c_p=0\) and
\(c_{p+1}\ne0\).  Note that \(c_0=\rho(1)\).

With the cut along the negative real axis, \(\log z\) is the principal
branch, holomorphic on \(\mathbb C\setminus(-\infty,0]\), which
contains a neighbourhood of \(z=1\); there \(e^{\log z}=z\) and
\[
  \log z=(z-1)E(z),\qquad E \text{ holomorphic},\ E(1)=1 .
\]
So \(\log z\) has a simple zero at \(z=1\) and no other zero.  The
cut is needed only to make \(\varphi\) a well-defined (single-valued)
holomorphic function near \(z=1\).

\step{Step 1: the key function \(g\)}
Define the entire function
\[
  g(h)=\rho(e^h)-h\,\sigma(e^h)=\sum_{j=0}^k\bigl(\alpha_j-h\beta_j\bigr)e^{jh}.
\]
This is \(L[e^x;h]\) at \(x=0\).  Expanding
\(e^{jh}=\sum_q j^qh^q/q!\),
\[
  g(h)=\sum_{q\ge0}h^q\Bigl[\frac1{q!}\sum_jj^q\alpha_j
       -\frac1{(q-1)!}\sum_jj^{q-1}\beta_j\Bigr]
  =\sum_{q\ge0}c_qh^q,
\]
(the \(\beta\) term absent for \(q=0\)).  So the \(c_q\) are exactly
the Taylor coefficients of \(g\) at \(h=0\), and
\begin{equation}\label{eq:ordg}
\begin{aligned}
  \text{order }p&\iff c_0=\cdots=c_p=0,\ c_{p+1}\ne0\\
  &\iff g\text{ has a zero of multiplicity exactly }p+1\text{ at }h=0 .
\end{aligned}
\end{equation}

\step{Step 2: relation between \(g\) and \(\varphi\)}
For \(z\) near \(1\), put \(h=\log z\).  Then
\[
  g(\log z)=\rho(z)-\log z\cdot\sigma(z)=\log z\cdot\varphi(z),
  \qquad z\ne1 .
\]
Conversely, for \(h\) near \(0\) (\(z=e^h\) near \(1\)),
\begin{equation}\label{eq:gphi}
  g(h)=h\,\varphi(e^h),\qquad h\ne0 .
\end{equation}

\step{Step 3: order \(p\Rightarrow\) zero of multiplicity \(p\)}
Suppose the method has order \(p\ge1\).  By \eqref{eq:ordg},
\(g(h)=h^{p+1}G(h)\) with \(G\) holomorphic and \(G(0)=c_{p+1}\ne0\).
In particular \(\rho(1)=c_0=0\), so the singularity of
\(\rho(z)/\log z\) at \(z=1\) is removable and \(\varphi\) is
holomorphic near \(1\).  Then
\[
  \varphi(z)=\frac{g(\log z)}{\log z}
  =(\log z)^pG(\log z)=(z-1)^p\,E(z)^p\,G(\log z),
\]
and \(E(1)^pG(\log1)=G(0)\ne0\).  Hence \(\varphi\) has a zero of
multiplicity exactly \(p\) at \(z=1\), with
\(\varphi(z)=c_{p+1}(z-1)^p+O\bigl((z-1)^{p+1}\bigr)\).

\step{Step 4: zero of multiplicity \(p\Rightarrow\) order \(p\)}
Suppose \(\varphi\) is holomorphic near \(z=1\) with a zero of
multiplicity \(p\ge1\) there.  (A zero requires holomorphy; if
\(\rho(1)\ne0\) then \(\rho/\log z\) has a pole at \(1\) and \(\varphi\)
cannot have a zero.  So \(\rho(1)=0\).)  Write
\(\varphi(z)=(z-1)^p\Phi(z)\) with \(\Phi(1)\ne0\).  Since
\(e^h-1=h\,\tilde E(h)\), \(\tilde E(0)=1\), \eqref{eq:gphi} gives
\[
  g(h)=h\,(e^h-1)^p\,\Phi(e^h)=h^{p+1}\,\tilde E(h)^p\,\Phi(e^h),
\]
with \(\tilde E(0)^p\Phi(1)=\Phi(1)\ne0\).  So \(g\) has a zero of
multiplicity exactly \(p+1\) at \(h=0\), and by \eqref{eq:ordg} the
method has order \(p\).

\step{Conclusion}
Both directions are proved: the method has order \(p\) if and only if
\(\varphi\) has a zero of multiplicity \(p\) at \(z=1\).  Moreover
\(\varphi(z)=c_{p+1}(z-1)^p+\cdots\), so the leading coefficient of
\(\varphi\) at \(z=1\) is the error constant \(c_{p+1}\).  This is
the basis of the construction of methods of maximal order: given
\(\rho\), take \(\sigma\) to be the Taylor polynomial of
\(\rho(z)/\log z\) at \(z=1\).
\end{solution}

%---------------------------------------------------------------------
\subsection{A tabulated solution and Numerov's method}

\begin{qbox}
A table of the function \(y=f(x)\) is given.
\begin{center}
\renewcommand{\arraystretch}{1.2}
\begin{tabular}{@{}c|ccccccc@{}}
\toprule
\(x\)&4&5&6&7&8&9&10\\
\(y\)&0.15024&0.56563&1.54068&3.25434&5.51438&7.56171&8.22108\\
\bottomrule
\end{tabular}
\end{center}
It is known that \(f\) satisfies the differential equation
\[
  \frac{d^2y}{dx^2}+\Bigl(1-\frac4x-\frac{n(n+1)}{x^2}\Bigr)y=0,
\]
where \(n\) is a positive integer.
\begin{enumerate}[label=(\alph*)]
  \item Find \(n\).\mk{3}
  \item Compute \(f(12)\) using Numerov's method with step size
    \(1\).\mk{3}
\end{enumerate}
\end{qbox}

\begin{solution}
Write the equation as \(y''=F(x,y)\) with
\[
  F(x,y)=-g(x)\,y,\qquad
  g(x)=1-\frac4x-\frac{N}{x^2},\qquad N=n(n+1).
\]
Numerov's method (order \(4\)) is
\[
  y_{i+1}-2y_i+y_{i-1}=\frac{h^2}{12}\bigl(F_{i+1}+10F_i+F_{i-1}\bigr).
\]

\step{(a) Finding \(n\)}
The tabulated values satisfy Numerov's relation up to its
truncation error \(O(h^6)\) and rounding.  Here \(h=1\) and
\(F_j=-(1-\frac4{x_j})y_j+N\,\dfrac{y_j}{x_j^2}\), which is linear in
\(N\).  Writing \(\delta^2y_i=y_{i+1}-2y_i+y_{i-1}\), the relation at
\(x_i\) gives
\[
  N=\frac{12\,\delta^2y_i+P_i}{Q_i},\qquad
  P_i=\sum_{j}w_j\Bigl(1-\frac4{x_j}\Bigr)y_j,\quad
  Q_i=\sum_jw_j\frac{y_j}{x_j^2},
\]
where the sums run over \(j=i-1,i,i+1\) with weights
\(w_{i\pm1}=1\), \(w_i=10\).  For example, at \(x_i=7\):
\(\delta^2y=5.51438-2(3.25434)+1.54068=0.54638\), and
\begin{align*}
  P&=\tfrac13(1.54068)+10\cdot\tfrac37(3.25434)+\tfrac12(5.51438)=17.21792,\\
  Q&=\tfrac{1.54068}{36}+10\cdot\tfrac{3.25434}{49}+\tfrac{5.51438}{64}=0.79311,\\
  N&=\frac{12(0.54638)+17.21792}{0.79311}=\frac{23.77448}{0.79311}=29.976 .
\end{align*}
At all interior points:
\begin{center}
\renewcommand{\arraystretch}{1.2}
\begin{tabular}{@{}cccccc@{}}
\toprule
\(x_i\)&5&6&7&8&9\\
\midrule
\(\delta^2y_i\)&0.55966&0.73861&0.54638&\(-0.21271\)&\(-1.38796\)\\
\(N\)&30.027&29.993&29.976&29.974&29.988\\
\bottomrule
\end{tabular}
\end{center}
(The cruder central-difference relation
\(\delta^2y_i=-g(x_i)y_i\) gives \(29.2\)--\(30.1\), the same
conclusion.)  Hence \(n(n+1)=30\), and since \(n\) is a positive
integer,
\[
  \boxed{n=5}\qquad(n(n+1)=30;\ \text{the other root } n=-6 \text{ is excluded}).
\]

\step{(b) \(f(12)\) by Numerov with \(h=1\)}
Now \(g(x)=1-\dfrac4x-\dfrac{30}{x^2}\).  Since \(F\) is linear in
\(y\), the implicit Numerov formula can be solved for \(y_{i+1}\):
\[
  \Bigl(1+\frac{g_{i+1}}{12}\Bigr)y_{i+1}
  =2\Bigl(1-\frac{5g_i}{12}\Bigr)y_i-\Bigl(1+\frac{g_{i-1}}{12}\Bigr)y_{i-1}.
\]
The two starting values are \(y(9)=7.56171\) and \(y(10)=8.22108\) from
the table.  The needed values of \(g\):
\[
  g(9)=\frac{5}{27}=0.185185,\quad
  g(10)=0.3,\quad
  g(11)=\frac{47}{121}=0.388430,\quad
  g(12)=\frac{11}{24}=0.458333 .
\]

\textit{Step to \(x=11\):}
\[
  1.032369\,y_{11}=2(0.875)(8.22108)-1.015432\,(7.56171)
  =14.386890-7.678403=6.708487,
\]
\[
  y_{11}=f(11)\approx6.49815 .
\]

\textit{Step to \(x=12\):}
\[
  1.038194\,y_{12}=2\Bigl(1-\frac{5(0.388430)}{12}\Bigr)(6.49815)-1.025\,(8.22108)
  =1.676309(6.49815)-8.426607,
\]
\[
  1.038194\,y_{12}=10.892900-8.426607=2.466293,
  \qquad
  \boxed{f(12)\approx2.37556 .}
\]

\step{Comment}
The equation is the Coulomb wave equation with \(\eta=2\), \(L=5\);
to the accuracy of the table, \(f(x)\approx6.2296\,F_5(2,x)\), the
regular Coulomb wave function, whose value gives \(f(12)\approx2.3818\).
The Numerov error \(\approx0.006\) is consistent with the large step
\(h=1\) (the local error is \(-\frac1{240}h^6y^{(6)}\), and the solution
oscillates with wavelength of order \(2\pi\)).
\end{solution}

%---------------------------------------------------------------------
\subsection{Du Fort--Frankel scheme: stability by the matrix method}

\begin{qbox}
Investigate the stability by the matrix method of the Du Fort and
Frankel scheme for the equation
\[
  \pd ut=\pdd ux,\qquad 0<x<1,\ t>0,
\]
for fixed mesh length, given that \(u\) is known on the boundaries
\(x=0\) and \(x=1\) for \(t>0\).  (Hint: the Du Fort and Frankel scheme
is
\[
  U_m^{n+1}=\frac{2r}{1+2r}\bigl(U_{m+1}^n+U_{m-1}^n\bigr)
  +\frac{1-2r}{1+2r}U_m^{n-1},\qquad r=\frac k{h^2}.)
\]\mk{5}
\end{qbox}

\begin{solution}
Let \(Nh=1\); the unknowns are \(U_m^n\), \(m=1,\dots,N-1\).

\step{Step 1: matrix form}
With \(U^n=[U_1^n,\dots,U_{N-1}^n]^T\) and the known boundary values
moved to the right-hand side,
\[
  U^{n+1}=\frac{2r}{1+2r}\,AU^n+\frac{1-2r}{1+2r}\,U^{n-1}+f_n,
  \qquad
  A=\begin{bmatrix}
    0&1&&\\
    1&0&1&\\
    &\ddots&\ddots&\ddots\\
    &&1&0
  \end{bmatrix}_{(N-1)\times(N-1)},
\]
\(f_n=\frac{2r}{1+2r}\,[U_0^n,0,\dots,0,U_N^n]^T\).

\step{Step 2: reduction to a two-level scheme}
This is a three-level scheme.  Put \(V^n=U^{n-1}\) and
\(W^n=\begin{bmatrix}U^n\\ V^n\end{bmatrix}\).  Then
\[
  W^{n+1}=BW^n+\begin{bmatrix}f_n\\0\end{bmatrix},
  \qquad
  B=\begin{bmatrix}
    \dfrac{2r}{1+2r}A&\dfrac{1-2r}{1+2r}I\\[8pt]
    I&0
  \end{bmatrix}.
\]
An error \(E^0\) in the initial data propagates as \(E^n=B^nE^0\).  For
a fixed mesh the scheme is stable if \(\norm{B^n}\) stays bounded as
\(n\to\infty\); this holds if every eigenvalue of \(B\) has modulus
\(<1\) (then \(B^n\to0\)), and fails if one has modulus \(>1\).

\step{Step 3: eigenvalues of \(B\)}
The eigenvalues of the tridiagonal matrix \(A\) are the distinct numbers
\[
  \lambda_k=2\cos\frac{k\pi}N,\qquad k=1,\dots,N-1,
\]
with eigenvectors \(v_k=\bigl(\sin\frac{jk\pi}N\bigr)_{j=1}^{N-1}\).
These \(N-1\) independent vectors are also eigenvectors of \(I\)
(eigenvalue \(1\)) and of \(0\) (eigenvalue \(0\)).  So all four blocks
of \(B\) share the eigenvectors \(v_k\), and by the theorem on
eigenvalues of block matrices, the eigenvalues \(\lambda\) of \(B\) are
the eigenvalues of the \(2\times2\) matrices
\[
  \begin{bmatrix}
    \dfrac{2r\lambda_k}{1+2r}&\dfrac{1-2r}{1+2r}\\[8pt]
    1&0
  \end{bmatrix},\qquad k=1,\dots,N-1 .
\]
(Directly: \(B\begin{bmatrix}\alpha_1v_k\\ \alpha_2v_k\end{bmatrix}
=\lambda\begin{bmatrix}\alpha_1v_k\\ \alpha_2v_k\end{bmatrix}\) reduces
to this \(2\times2\) eigenproblem for \((\alpha_1,\alpha_2)\).)  Hence
\[
  \lambda^2-\frac{2r\lambda_k}{1+2r}\lambda-\frac{1-2r}{1+2r}=0,
  \qquad
  \lambda=\frac{r\lambda_k\pm\bigl(1+(\lambda_k^2-4)r^2\bigr)^{1/2}}{1+2r},
\]
and with \(\lambda_k=2\cos\frac{k\pi}N\), \(\lambda_k^2-4=-4\sin^2\frac{k\pi}N\):
\[
  \lambda=\frac{2r\cos\frac{k\pi}N\pm\bigl(1-4r^2\sin^2\frac{k\pi}N\bigr)^{1/2}}{1+2r},
  \qquad k=1,\dots,N-1 .
\]
For \(1\le k\le N-1\), \(\sin\frac{k\pi}N>0\) and
\(|\cos\frac{k\pi}N|<1\).

\step{Case I: \(0\le1-4r^2\sin^2\frac{k\pi}N<1\) (real roots)}
\[
  |\lambda|\le\frac{2r|\cos\frac{k\pi}N|+\bigl(1-4r^2\sin^2\frac{k\pi}N\bigr)^{1/2}}{1+2r}
  <\frac{2r+1}{1+2r}=1 .
\]

\step{Case II: \(1-4r^2\sin^2\frac{k\pi}N<0\) (complex conjugate roots)}
This requires \(4r^2>1\), i.e.\ \(r>\frac12\).  Then
\[
  |\lambda|^2
  =\frac{4r^2\cos^2\frac{k\pi}N+4r^2\sin^2\frac{k\pi}N-1}{(1+2r)^2}
  =\frac{4r^2-1}{(1+2r)^2}=\frac{2r-1}{2r+1}<1 .
\]
(Equivalently, \(|\lambda|^2\) is the product of the roots,
\(-\frac{1-2r}{1+2r}\).)

\step{Conclusion}
Every eigenvalue of \(B\) satisfies \(|\lambda|<1\) for every \(r>0\).
Hence, for a fixed mesh size, \(B^n\to0\) and the Du Fort--Frankel
scheme is \emph{unconditionally stable} (for all \(r=k/h^2>0\)).

\textit{Remark.}  Stability is unconditional, but the scheme is
consistent with \(u_t=u_{xx}\) only if \(k/h\to0\); if \(k/h=c\) is
fixed, it is consistent with \(u_t=u_{xx}-c^2u_{tt}\) instead.
\end{solution}

%---------------------------------------------------------------------
\subsection{Wave equation by an implicit three-level scheme}

\begin{qbox}
Find the solution of the initial boundary value problem
\[
  \pdd ut=\pdd ux,\qquad 0<x<1,
\]
subject to the initial conditions \(u(x,0)=\sin\pi x\),
\(\pd ut(x,0)=0\), \(0\le x\le1\), and the boundary conditions
\(u(0,t)=0\), \(u(1,t)=0\), \(t>0\), by using the implicit scheme
\[
  \delta_t^2U_m^n=r\,\delta_x^2\bigl\{\theta U_m^{n+1}+(1-2\theta)U_m^n
  +\theta U_m^{n-1}\bigr\}\quad\text{with }\theta=\tfrac12,
\]
where \(r=\frac kh\).  Assume \(h=\frac14\), \(r=\frac34\) and
investigate for one time step.\mk{5}
\end{qbox}

\begin{solution}
We use the scheme exactly as printed, with
\[
  r=\frac kh=\frac34,\qquad h=\frac14,\qquad k=rh=\frac3{16},\qquad \theta=\frac12 .
\]

\step{Step 1: the scheme for \(\theta=\frac12\)}
With \(\theta=\frac12\), \(1-2\theta=0\), so
\[
  U_m^{n+1}-2U_m^n+U_m^{n-1}
  =\frac r2\,\delta_x^2\bigl(U_m^{n+1}+U_m^{n-1}\bigr).
\]
Since \(\theta=\frac12\ge\frac14\), the scheme is unconditionally
stable (von Neumann: the amplification factors satisfy
\(|\xi_1|=|\xi_2|=1\) for every value of the positive coefficient of
\(\delta_x^2\)).

\step{Step 2: the first time step}
Initial values: \(U_m^0=\sin\pi mh\).  The condition \(u_t(x,0)=0\) is
replaced by the central difference
\[
  \frac{U_m^1-U_m^{-1}}{2k}=0\quad\Longrightarrow\quad U_m^{-1}=U_m^1 .
\]
Putting \(n=0\) in the scheme and eliminating \(U_m^{-1}\):
\[
  2U_m^1-2U_m^0=\frac r2\,\delta_x^2\bigl(2U_m^1\bigr)
  \quad\Longrightarrow\quad
  \Bigl(1-\frac r2\delta_x^2\Bigr)U_m^1=U_m^0,
\]
that is,
\[
  (1+r)U_m^1-\frac r2\bigl(U_{m+1}^1+U_{m-1}^1\bigr)=U_m^0 .
\]
With \(r=\frac34\), multiplying by \(8\):
\[
  14\,U_m^1-3\bigl(U_{m+1}^1+U_{m-1}^1\bigr)=8\,U_m^0,
  \qquad m=1,2,3,
\]
with \(U_0^1=U_4^1=0\).

\step{Step 3: solving}
\(h=\frac14\), so the interior points are \(x=\frac14,\frac12,\frac34\)
with \(U^0=\bigl(\frac1{\sqrt2},1,\frac1{\sqrt2}\bigr)\).  The problem is
symmetric about \(x=\frac12\), so \(U_1^1=U_3^1=a\), \(U_2^1=b\):
\begin{align*}
  m=1:&\quad 14a-3b=\frac{8}{\sqrt2}=4\sqrt2,\\
  m=2:&\quad 14b-6a=8 .
\end{align*}
From the first, \(b=(14a-4\sqrt2)/3\); substituting,
\(196a-56\sqrt2-18a=24\), so
\[
  a=\frac{12+28\sqrt2}{89}\approx0.57975,\qquad
  b=\frac{56+12\sqrt2}{89}\approx0.81989 .
\]
Hence, at \(t=k=\frac3{16}\),
\[
  \boxed{U^1\bigl(\tfrac14\bigr)=U^1\bigl(\tfrac34\bigr)\approx0.57975,\qquad
  U^1\bigl(\tfrac12\bigr)\approx0.81989.}
\]
(Check: \(14(0.57975)-3(0.81989)=5.65684\approx4\sqrt2\) and
\(14(0.81989)-6(0.57975)=8.0000\).)

\step{Step 4: investigation (comparison with the exact solution)}
The exact solution (by separation of variables) is
\(u(x,t)=\sin\pi x\cos\pi t\).  Note that
\(\delta_x^2\sin\pi mh=-4\sin^2\frac{\pi h}2\,\sin\pi mh=-(2-\sqrt2)\sin\pi mh\)
for \(h=\frac14\): the initial data is an eigenvector of
\(\delta_x^2\), so the numerical solution keeps the exact shape
\(\sin\pi x\), and only its amplitude is approximated:
\[
  U_m^1=\frac{\sin\pi mh}{1+2r\sin^2\frac{\pi h}2}
  =\frac{\sin\pi mh}{1+\frac38(2-\sqrt2)}=0.81989\,\sin\pi mh,
\]
in agreement with Step 3.  The exact amplitude is
\(\cos\frac{3\pi}{16}=0.83147\).
\begin{center}
\renewcommand{\arraystretch}{1.2}
\begin{tabular}{@{}cccc@{}}
\toprule
\(x\)&numerical \(U^1\)&exact \(u(x,\frac3{16})\)&error\\
\midrule
\(\frac14,\frac34\)&0.57975&0.58794&0.0082\\
\(\frac12\)&0.81989&0.83147&0.0116\\
\bottomrule
\end{tabular}
\end{center}
The relative error is about \(1.4\%\) on this coarse grid.

\textit{Remark.}  Deriving the scheme from
\(\frac1{k^2}\delta_t^2U=\frac1{h^2}\delta_x^2U\) gives the coefficient
\((k/h)^2=r^2\) (as in the notes).  With \(r^2=\frac9{16}\) in place of
\(r\), the same steps give
\(50U_m^1-9(U_{m+1}^1+U_{m-1}^1)=32U_m^0\) and
\(U^1=0.60709,\ 0.85855,\ 0.60709\).
\end{solution}

%---------------------------------------------------------------------
\subsection{Poisson's equation on a square: five-point scheme}

\begin{qbox}
Consider the Poisson equation with Dirichlet boundary conditions:
\[
  \pdd ux+\pdd uy=16,\qquad |x|<1,\ |y|<1;\qquad u=0\ \text{on the boundary } |x|=1,\ |y|=1.
\]
Formulate the five-point difference scheme with mesh size \(h\) in both
directions.  Solve the difference equations for \(h=\frac12\).\mk{5}
\end{qbox}

\begin{solution}
\step{Step 1: the five-point scheme}
Let \(Nh=2\), \(x_i=-1+ih\), \(y_j=-1+jh\), \(0\le i,j\le N\), and
\(u_{i,j}\approx u(x_i,y_j)\).  Using
\(\pdd ux=\frac1{h^2}\delta_x^2u+O(h^2)\) and similarly in \(y\), at each
interior point
\[
  \frac{u_{i+1,j}+u_{i-1,j}+u_{i,j+1}+u_{i,j-1}-4u_{i,j}}{h^2}=16,
  \qquad 1\le i,j\le N-1,
\]
i.e.
\[
  \boxed{\;u_{i+1,j}+u_{i-1,j}+u_{i,j+1}+u_{i,j-1}-4u_{i,j}=16h^2,\;}
\]
with \(u_{i,j}=0\) whenever \(i\in\{0,N\}\) or \(j\in\{0,N\}\).  The
local truncation error is \(O(h^2)\).  In matrix form, ordering the
unknowns row by row,
\(\operatorname{tridiag}(I,T,I)\,\mathbf u=16h^2\mathbf 1\) with
\(T=\operatorname{tridiag}(1,-4,1)\), each block of order \(N-1\).  The
matrix is symmetric and negative definite, so the solution is unique.

\step{Step 2: \(h=\frac12\)}
Then \(N=4\), the grid lines are \(-1,-\frac12,0,\frac12,1\), there are
\(3\times3=9\) interior points, and \(16h^2=4\).  The problem is
invariant under \(x\mapsto-x\), \(y\mapsto-y\) and \(x\leftrightarrow y\),
and so is the difference system; by uniqueness its solution has the
same symmetry.  So there are only three distinct unknowns:
\[
  \alpha=u\bigl(\pm\tfrac12,\pm\tfrac12\bigr),\qquad
  \beta=u\bigl(0,\pm\tfrac12\bigr)=u\bigl(\pm\tfrac12,0\bigr),\qquad
  \gamma=u(0,0).
\]
\begin{itemize}
  \item At \((\frac12,\frac12)\): neighbours \((1,\frac12)\),
    \((\frac12,1)\) on the boundary, \((0,\frac12)\), \((\frac12,0)\):
    \[ 2\beta-4\alpha=4 . \]
  \item At \((0,\frac12)\): neighbours \((\pm\frac12,\frac12)\),
    \((0,1)\) on the boundary, \((0,0)\):
    \[ 2\alpha+\gamma-4\beta=4 . \]
  \item At \((0,0)\): four neighbours of type \(\beta\):
    \[ 4\beta-4\gamma=4 . \]
\end{itemize}
From the first and third, \(\alpha=\frac\beta2-1\) and
\(\gamma=\beta-1\).  Substituting in the second,
\(\beta-2+\beta-1-4\beta=4\), so \(\beta=-\frac72\).  Hence
\[
  \boxed{\alpha=u\bigl(\pm\tfrac12,\pm\tfrac12\bigr)=-\frac{11}4,\qquad
  \beta=u\bigl(0,\pm\tfrac12\bigr)=u\bigl(\pm\tfrac12,0\bigr)=-\frac72,\qquad
  \gamma=u(0,0)=-\frac92 .}
\]

\step{Check}
\(2(-\frac72)-4(-\frac{11}4)=4\);
\(2(-\frac{11}4)-\frac92+14=4\);
\(4(-\frac72)+18=4\).  On the grid,
\[
  \begin{array}{c|ccc}
    y\backslash x&-\frac12&0&\frac12\\\hline
    \frac12&-2.75&-3.5&-2.75\\
    0&-3.5&-4.5&-3.5\\
    -\frac12&-2.75&-3.5&-2.75
  \end{array}
\]
The solution is negative, as it must be (\(\Delta u>0\) with zero
boundary values, by the maximum principle).  For comparison, the exact
value (from the Fourier series solution) is \(u(0,0)\approx-4.715\).
\end{solution}

%---------------------------------------------------------------------
\subsection{Peaceman--Rachford ADI for the Neumann problem}

\begin{qbox}
Consider Laplace's equation
\(\pdd ux+\pdd uy=0\), \(0<x,y<1\), subject to \(\pd un=g(x,y)\),
where \(\pd un\) is the outward directed normal derivative on the
boundary of the unit square \(0\le x,y\le1\).  The square region is
covered by a grid with sides parallel to the coordinate axes and the
grid spacing is \(h\) and \(Mh=1\), where \(M\) is the number of
intervals in each direction.  Prove that the Peaceman--Rachford
Alternating Direction Implicit (A.D.I.) iterative method for solving
Laplace's equation in a square with Neumann boundary conditions is
convergent for any fixed value of the iteration parameter.\mk{5}
\end{qbox}

\begin{solution}
\step{Step 1: the discrete problem}
Since \(u\) is not prescribed on the boundary, all \((M+1)^2\) grid
values \(u_{l,m}\), \(0\le l,m\le M\), are unknowns.  At every grid
point use the five-point formula
\[
  4u_{l,m}-u_{l+1,m}-u_{l-1,m}-u_{l,m+1}-u_{l,m-1}=0 .
\]
At boundary points it involves fictitious values outside the square,
which are eliminated with central differences for the normal
derivative.  For instance on \(x=0\), \(-\frac{u_{1,m}-u_{-1,m}}{2h}=g_{0,m}\),
so \(u_{-1,m}=u_{1,m}+2hg_{0,m}\), giving
\[
  4u_{0,m}-2u_{1,m}-u_{0,m+1}-u_{0,m-1}=2hg_{0,m},
\]
and at the corner \((0,0)\),
\(4u_{0,0}-2u_{1,0}-2u_{0,1}=2h\cdot2g_{0,0}\); the other sides and
corners are similar.  With \(U\) ordered row by row, this gives
\[
  AU=2hG,
\]
where \(G\) has entry \(2g\) at the corners, \(g\) at the other boundary
points and \(0\) at interior points.

\step{Step 2: the splitting \(A=X+Y\)}
Let
\[
  C=\begin{bmatrix}
    2&-2&&&\\
    -1&2&-1&&\\
    &\ddots&\ddots&\ddots&\\
    &&-1&2&-1\\
    &&&-2&2
  \end{bmatrix}_{(M+1)\times(M+1)}
\]
represent the \(x\)-difference \(2u_{l,m}-u_{l+1,m}-u_{l-1,m}\)
(including the boundary rows).  Then \(A=X+Y\) with
\[
  X=I\otimes C=\operatorname{diag}(C,\dots,C),\qquad
  Y=C\otimes I=\begin{bmatrix}
    2I&-2I&&\\
    -I&2I&-I&\\
    &\ddots&\ddots&\ddots\\
    &&-2I&2I
  \end{bmatrix},
\]
the \(x\)- and \(y\)-difference parts.  (Indeed
\(X+Y=\operatorname{tridiag}(-I,\ C+2I,\ -I)\) with first and last
off-diagonal blocks \(-2I\), which is \(A\).)

\step{Step 3: eigenvectors of \(C\), \(X\), \(Y\)}
Let \(c_p=\bigl(\cos\frac{lp\pi}M\bigr)_{l=0}^M\),
\(p=0,1,\dots,M\).  For \(0<l<M\),
\(-\cos\frac{(l-1)p\pi}M+2\cos\frac{lp\pi}M-\cos\frac{(l+1)p\pi}M
=\bigl(2-2\cos\frac{p\pi}M\bigr)\cos\frac{lp\pi}M\).  For \(l=0\):
\(2-2\cos\frac{p\pi}M\).  For \(l=M\):
\(-2\cos\frac{(M-1)p\pi}M+2\cos p\pi=(2-2\cos\frac{p\pi}M)\cos p\pi\),
using \(\cos\frac{(M-1)p\pi}M=\cos p\pi\cos\frac{p\pi}M\).  Hence
\[
  Cc_p=\mu_pc_p,\qquad \mu_p=4\sin^2\frac{p\pi}{2M},\qquad p=0,\dots,M .
\]
The \(\mu_p\) are \(M+1\) distinct numbers in \([0,4]\), so the \(c_p\)
form a basis.  Consequently the \((M+1)^2\) vectors
\(\beta_{p,q}=c_q\otimes c_p\) (components
\(\cos\frac{lp\pi}M\cos\frac{mq\pi}M\)) form a basis, and
\[
  X\beta_{p,q}=\mu_p\beta_{p,q},\qquad
  Y\beta_{p,q}=\mu_q\beta_{p,q},\qquad
  A\beta_{p,q}=(\mu_p+\mu_q)\beta_{p,q}.
\]
In particular \(XY=YX\), \(\mu_p\ge0\), and \(\mu_p=0\iff p=0\).  So
\(A\) is singular with null space spanned by
\(\beta_{0,0}=\mathbf 1\) (the constant vector), and the range of \(A\)
is spanned by the \(\beta_{p,q}\), \((p,q)\ne(0,0)\).  The discrete
Neumann problem is solvable if and only if the \(\beta_{0,0}\)-component
of \(G\) vanishes (the discrete compatibility condition, see the remark
below), and then its solution is unique up to an additive constant.
We assume this, as for the continuous problem.

\step{Step 4: the Peaceman--Rachford iteration}
With a fixed iteration parameter \(\rho>0\) (in the notes' notation
\(\rho=2/\lambda\)),
\begin{align*}
  (\rho I+X)\,U^{n+\frac12}&=(\rho I-Y)\,U^n+2hG,\\
  (\rho I+Y)\,U^{n+1}&=(\rho I-X)\,U^{n+\frac12}+2hG .
\end{align*}
The matrices \(\rho I+X\) and \(\rho I+Y\) have eigenvalues
\(\rho+\mu_p>0\), so they are invertible (and each is a set of
tridiagonal systems along grid lines).  Any solution \(U\) of
\(AU=2hG\) is a fixed point: \((\rho I+X)U=(\rho I-Y)U+(X+Y)U=(\rho I-Y)U+2hG\),
and similarly for the second half-step.  Hence the error
\(e^n=U^n-U\) satisfies
\[
  e^{n+1}=T\,e^n,\qquad
  T=(\rho I+Y)^{-1}(\rho I-X)(\rho I+X)^{-1}(\rho I-Y).
\]

\step{Step 5: eigenvalues of the iteration matrix}
Because \(X\) and \(Y\) share the eigenvectors \(\beta_{p,q}\),
\[
  T\beta_{p,q}=\tau_{p,q}\beta_{p,q},\qquad
  \tau_{p,q}=\frac{(\rho-\mu_p)(\rho-\mu_q)}{(\rho+\mu_p)(\rho+\mu_q)} .
\]
For \(\rho>0\) and \(\mu\ge0\), \(\bigl|\frac{\rho-\mu}{\rho+\mu}\bigr|\le1\),
with equality only for \(\mu=0\).  Therefore
\[
  \tau_{0,0}=1,\qquad |\tau_{p,q}|<1\quad\text{for all }(p,q)\ne(0,0),
\]
for \emph{every} \(\rho>0\).

\step{Step 6: convergence}
Expand \(e^0=\sum_{p,q}\varepsilon_{p,q}\beta_{p,q}\).  Then
\[
  e^n=\varepsilon_{0,0}\mathbf 1+\sum_{(p,q)\ne(0,0)}\tau_{p,q}^n\,\varepsilon_{p,q}\beta_{p,q}
  \ \longrightarrow\ \varepsilon_{0,0}\mathbf 1\qquad(n\to\infty),
\]
geometrically, with factor \(\max_{(p,q)\ne(0,0)}|\tau_{p,q}|<1\).
Hence
\[
  U^n\to U+\varepsilon_{0,0}\mathbf 1,
\]
and since \(A\mathbf 1=0\), the limit is again a solution of
\(AU=2hG\).  The only component that does not decay is the constant
vector, which is exactly the non-uniqueness of the Neumann problem
(the constant is fixed by the initial guess).  So the
Peaceman--Rachford iteration converges to a solution of the discrete
Neumann problem for any fixed value of the iteration parameter.
\(\blacksquare\)

\step{Remark: the compatibility condition}
The vector \(w=(\frac12,1,\dots,1,\frac12)^T\) satisfies \(w^TC=0\), so
\((w\otimes w)^TA=0\), and \(AU=2hG\) is solvable if and only if
\((w\otimes w)^TG=0\).  The weights \(w_lw_m\) are \(\frac14\) at the
corners (where \(G=2g\)) and \(\frac12\) at the other boundary points
(where \(G=g\)), so the condition is \(\sum_{\text{boundary points}}g=0\):
the trapezoidal-rule form of \(\oint_{\partial R}g\,ds=0\), which is the
condition for the continuous Neumann problem to be solvable.  If it
fails, the \(\beta_{0,0}\)-component of the iterates drifts linearly
while all other components still converge.
\end{solution}

\end{document}
